% =============================================================================
%  Analytical T-Matrix for Elastic Scatterers:
%  Sphere and Cube Geometries
%  Derived from the Lippmann--Schwinger Integral Equation
%
%  Formatted for Geophysical Journal International
% =============================================================================
\documentclass[extra,mreferee]{gji}

% --- Packages ----------------------------------------------------------------
\usepackage{timet}              % GJI Times font package
\usepackage{amsmath,amssymb,amsfonts}
\usepackage{mathtools}
\usepackage{bm}
\usepackage{array}
\usepackage{booktabs}
\usepackage{graphicx}
\usepackage{url}

% --- Custom commands ---------------------------------------------------------
\newcommand{\bx}{\mathbf{x}}
\newcommand{\bG}{\mathbf{G}}
\newcommand{\bT}{\mathbf{T}}
\newcommand{\bI}{\mathbf{I}}
\newcommand{\bJ}{\mathbf{J}}
\newcommand{\bD}{\mathbf{D}}
\newcommand{\bS}{\mathbf{S}}
\newcommand{\bP}{\mathbf{P}}
\newcommand{\bv}{\mathbf{v}}
\newcommand{\bu}{\mathbf{u}}
\newcommand{\bt}{\mathbf{t}}
\newcommand{\Hmat}{\mathcal{H}}

% --- Title and authors -------------------------------------------------------
\title[T-Matrix for Elastic Scatterers]{%
  Analytical T-Matrix for Elastic Scatterers:
  Sphere and Cube Geometries ---
  Derived from the Lippmann--Schwinger Integral Equation}

\author[T.\,M.\ Nestor]{%
  T.\,M.\ Nestor$^{1}$ \\
  $^{1}$Research School of Earth Sciences,
  Australian National University, Canberra, ACT, Australia
}

\date{Received \today; in original form \today}

% =============================================================================
\begin{document}

\maketitle

% --- Summary -----------------------------------------------------------------
\begin{summary}
We apply symbolic computer algebra (SymPy) to evaluate the
Lippmann--Schwinger integral equation analytically for two scatterer
geometries: a sphere and a cube, both with constant isotropic contrasts
$\Delta\lambda$, $\Delta\mu$, $\Delta\rho$, embedded in a homogeneous
isotropic background.  The sphere yields closed-form results with isotropic
effective contrasts; the cube introduces cubic anisotropy in the effective
stiffness, providing the natural building block for effective medium
theories of heterogeneous solids.  Self-consistent amplification factors
are derived that sum the multiple-scattering series inside each inclusion
exactly, yielding a T-matrix that satisfies the optical theorem.  The local
site T-matrix is then embedded in the Foldy--Lax multiple-scattering
framework with the inter-site Green's tensor provided by the Kennett
reflectivity propagator for layered media.  A hybrid numerical
implementation in NumPy achieves $\sim\!10^{-8}$ relative accuracy by
combining Gauss--Legendre quadrature for the convergent Green's tensor
integral with a polynomial Taylor expansion that avoids the divergent
Eshelby singularity in the second-derivative integrals.  The displacement-traction
$6\times 6$ T-matrix is derived for direct integration into the Kennett
propagator framework.
\end{summary}

% --- Keywords ----------------------------------------------------------------
\begin{keywords}
  Seismic wave propagation --
  Computational seismology --
  Wave scattering and diffraction --
  Numerical modelling --
  Theoretical seismology.
\end{keywords}

% --- Body chapters -----------------------------------------------------------
% =============================================================================
%  Part I: Spherical Scatterer  (Sections 1--10)
% =============================================================================

\section{Introduction}

The scattering of elastic waves by localised heterogeneities is a fundamental
problem in seismology, ultrasonics, and geophysical exploration.  When the
scatterer is a sphere of radius~$a$ embedded in a homogeneous isotropic
background, the problem admits an exact solution in terms of partial-wave (Mie)
series \citep{Korneev1993,Gritto1995}.  In the long-wavelength (Rayleigh) regime
$ka\ll 1$, only the lowest-order multipoles contribute and the scattered field
can be expressed compactly through an effective T-matrix that maps the incident
wavefield to the scattered wavefield.

The T-matrix concept originates in quantum scattering theory and was adapted to
elastic-wave problems by \citet{Waterman1976}; \citet{Bostock1991}
and \citet{BostockKennett1992} applied it to surface-wave
scattering.  \citet{Pollitz1998} derived the scattered field for a small
spherical inclusion without assuming small perturbations, arriving at
self-consistent amplification factors that are functionally equivalent to
renormalised material contrasts.  \citet{Pellegrini2011}
formalised this as a mean-field T-matrix approach: by approximating the interior
stress and momentum fields by their volume averages, the full
multiple-scattering series inside the inclusion can be summed in closed form,
yielding a T-matrix that satisfies the optical theorem exactly.

On the computational side, Jakobsen and co-workers developed integral-equation
methods for seismic modelling and inversion based on the Lippmann--Schwinger
equation.  \citet{Jakobsen2012} introduced the T-matrix formalism for
acoustic forward modelling; \citet{JakobsenUrsin2015} extended
this to direct iterative T-matrix inversion; and \citet{JakobsenWu2016}
proposed a renormalised scattering series for frequency-domain waveform modelling
in the presence of strong velocity contrasts, overcoming the convergence
limitations of the Born series.  Most recently, \citet{Shekhar2024} formulated a numerical integral-equation method for
microseismic wavefield modelling in arbitrarily anisotropic elastic media,
providing the Green's tensor expressions and the Lippmann--Schwinger framework
used as the starting point of this work.

In this note, we apply symbolic computer algebra (SymPy) to evaluate the
Lippmann--Schwinger integral equation analytically for two scatterer geometries:
a sphere and a cube, both with constant isotropic contrasts $\Delta\lambda$,
$\Delta\mu$, $\Delta\rho$.  The sphere yields closed-form results with
isotropic effective contrasts; the cube introduces cubic anisotropy in the
effective stiffness, providing the natural building block for effective medium
theories of heterogeneous solids.


% =====================================================================
\part*{Part I: Spherical Scatterer}
% =====================================================================

\section{Problem Statement}

We consider the scattering of elastic waves by a small spherical inclusion of
radius~$a$ embedded in an infinite homogeneous isotropic background medium with
P-wave velocity~$\alpha$, S-wave velocity~$\beta$, and density~$\rho$.  The
inclusion has constant isotropic contrasts $\Delta\lambda$, $\Delta\mu$,
$\Delta\rho$ relative to the background.

The wavefield satisfies the Lippmann--Schwinger integral equation (Eq.~8
of \citealt{Shekhar2024}):
\begin{equation}
  u_i(\bx) = u_i^0(\bx) + \int_D \Bigl[
    G_{ij}(\bx,\bx')\,\omega^2\,\Delta\rho\;u_j(\bx')
    + H_{ijk}(\bx,\bx')\,\Delta c_{jklm}\,\varepsilon_{lm}(\bx')
  \Bigr] d^3x',
\label{eq:LS}
\end{equation}
where the strain--Green's tensor is
\begin{equation}
  H_{ijk}^{(0)}(\bx,\bx')
  = \tfrac{1}{2}\bigl(G_{ij,k} + G_{ik,j}\bigr),
\label{eq:Hdef}
\end{equation}
and the isotropic stiffness contrast is
\begin{equation}
  \Delta c_{jklm}
  = \Delta\lambda\,\delta_{jk}\,\delta_{lm}
  + \Delta\mu\,(\delta_{jl}\,\delta_{km} + \delta_{jm}\,\delta_{kl}).
\label{eq:Dc}
\end{equation}


\section{Background Green's Tensor}

The elastodynamic Green's tensor for a homogeneous isotropic medium
(Appendix~A of \citealt{Shekhar2024}) admits the radial decomposition:
\begin{equation}
  G_{ij}(\bx) = f(r)\,\delta_{ij} + g(r)\,\hat{n}_i\,\hat{n}_j,
\label{eq:Grad}
\end{equation}
with $\hat{n}_i = x_i/r$ and radial functions
\begin{equation}
  X(r) = \frac{e^{i\omega r/\alpha}}{4\pi\rho\alpha^2}, \quad
  V(r) = \frac{e^{i\omega r/\beta}}{4\pi\rho\beta^2}, \quad
  C = \frac{1-\beta^2/\alpha^2}{8\pi\rho\beta^2},
  \qquad
  f(r) = \frac{V(r)}{r} - \frac{C}{r},
\label{eq:fdef}
\end{equation}
\begin{equation}
  g(r) = \frac{3C}{r} - \frac{X(r)}{r} + \frac{V(r)}{r} .
\label{eq:gdef}
\end{equation}


\section{Calculus: Spherical Coordinate Reduction}

The central computational challenge is to evaluate volume integrals of the
form $\int_{\text{sphere}} G_{ij}\,d^3x$,
$\int_{\text{sphere}} G_{ij,k}\,d^3x$ and
$\int_{\text{sphere}} G_{ij,kl}\,d^3x$ over a sphere of radius~$a$ centred
at the scatterer.  The key insight is that although these integrals live in
Cartesian tensor space, the radial decomposition~\eqref{eq:Grad} allows them
to be factored into angular integrals (which evaluate to known isotropic
tensor identities) multiplied by 1D radial integrals (which SymPy evaluates
in closed form).

\subsection{Conversion to Spherical Coordinates}

We write the volume element as $d^3x = r^2\,dr\,d\Omega$, where
$d\Omega = \sin\theta\,d\theta\,d\varphi$ is the solid-angle element.  Every
Cartesian integrand is expressed in terms of the radial coordinate~$r$ and
the unit direction vector $\hat{n}_i = x_i/r$.  Because $f(r)$ and $g(r)$
depend only on~$r$, the angular and radial parts separate cleanly:
\begin{equation}
  \int_{\text{sphere}}
    \bigl[f(r)\,\delta_{ij} + g(r)\,\hat{n}_i\,\hat{n}_j\bigr]
    r^2\,dr\,d\Omega
  = \delta_{ij}\int_0^a r^2 f(r)\,dr \int d\Omega
  + \int_0^a r^2 g(r)\,dr \int d\Omega\;\hat{n}_i\hat{n}_j.
\label{eq:sphfact}
\end{equation}

\subsection{Angular Integral Identities}

The angular integrals of products of direction cosines over the unit sphere
are determined entirely by symmetry.  The only non-vanishing cases with an
even number of $\hat{n}$ factors are:
\begin{equation}
  \int d\Omega = 4\pi,
\label{eq:ang0}
\end{equation}
\begin{equation}
  \int d\Omega\;\hat{n}_i\,\hat{n}_j = \frac{4\pi}{3}\,\delta_{ij},
\label{eq:ang2}
\end{equation}
\begin{equation}
  \int d\Omega\;\hat{n}_i\,\hat{n}_j\,\hat{n}_k\,\hat{n}_l
  = \frac{4\pi}{15}\,(\delta_{ij}\delta_{kl}
    + \delta_{ik}\delta_{jl} + \delta_{il}\delta_{jk}).
\label{eq:ang4}
\end{equation}
Integrals with an odd number of $\hat{n}$ factors vanish by parity.  This is
the reason that $\int H_{ijk}\,d^3x = 0$: each term in $G_{ij,k}$ contains
an odd number of direction cosines after the chain rule is applied, so the
angular integral is zero.

\subsection{Chain-Rule Decomposition of Derivatives}

The Cartesian derivatives of $G_{ij}$ must be expressed in terms of
$\hat{n}_i$ and radial derivatives of $f$ and~$g$ before integration.  We
use the fundamental identities:
\begin{equation}
  \frac{\partial r}{\partial x_k} = \hat{n}_k, \qquad
  \frac{\partial\hat{n}_i}{\partial x_k}
  = \frac{1}{r}\bigl(\delta_{ik} - \hat{n}_i\,\hat{n}_k\bigr).
\label{eq:chain}
\end{equation}

\textbf{First derivative.}  Applying the chain rule to
$G_{ij} = f\,\delta_{ij} + g\,\hat{n}_i\hat{n}_j$:
\begin{equation}
  G_{ij,k} = f'\,\delta_{ij}\,\hat{n}_k
  + g'\,\hat{n}_i\,\hat{n}_j\,\hat{n}_k
  + \frac{g}{r}\bigl(\delta_{ik}\,\hat{n}_j
    + \delta_{jk}\,\hat{n}_i - 2\hat{n}_i\,\hat{n}_j\,\hat{n}_k\bigr),
\label{eq:Gijk}
\end{equation}
where primes denote $d/dr$.  Every term contains an odd number of
$\hat{n}$ factors (one or three), confirming that
$\int G_{ij,k}\,d^3x = 0$.

\textbf{Second derivative.}  Differentiating~\eqref{eq:Gijk} with respect
to~$x_l$ and again applying~\eqref{eq:chain} produces terms with angular
structures involving 0, 2, or 4 direction cosines:
\begin{equation}
  G_{ij,kl} = \text{terms proportional to}
  \;\delta\delta,\;\delta\hat{n}\hat{n},\;\hat{n}\hat{n}\hat{n}\hat{n},
\label{eq:Gijkl}
\end{equation}
with coefficients built from $f$, $g$ and their first and second radial
derivatives.  Rather than writing out all terms explicitly, we use a
\emph{contraction strategy}: the full tensor
$\int G_{ij,kl}\,d^3x$ must have the isotropic form
$A\,\delta_{ij}\delta_{kl} + B\,(\delta_{ik}\delta_{jl}+\delta_{il}\delta_{jk})$,
and we can determine the two unknowns $A$ and~$B$ from any two independent
contractions.

\subsection{Contraction Strategy for $A$ and $B$}

\textbf{Contraction~1:}  Trace over $(i,j)$.  Setting $i=j$ and summing
gives $G_{ii} = 3f + g \equiv h(r)$, which is a scalar function of~$r$.
Its second Cartesian derivative is simply:
\begin{equation}
  G_{ii,kl} = h''(r)\,\hat{n}_k\,\hat{n}_l
  + \frac{h'(r)}{r}\bigl(\delta_{kl} - \hat{n}_k\,\hat{n}_l\bigr).
\label{eq:trace1}
\end{equation}
Integrating over the sphere using \eqref{eq:ang0}--\eqref{eq:ang2}:
\begin{equation}
  \int G_{ii,kl}\,d^3x
  = \delta_{kl} \int_0^a r^2
    \left[\frac{4\pi}{3}\,h'' + \frac{8\pi}{3}\,\frac{h'}{r}\right] dr
  = (3A + 2B)\,\delta_{kl}.
\label{eq:contr1}
\end{equation}

\textbf{Contraction~2:}  Divergence structure.  Summing over $j=k$ gives
the divergence $G_{ij,j} = q(r)\,\hat{n}_i$, where
$q(r) = f' + g' + 2g/r$.  Its further derivative with respect to~$x_l$ has
the same scalar-times-vector structure:
\begin{equation}
  G_{ij,jl} = q'(r)\,\hat{n}_i\,\hat{n}_l
  + \frac{q(r)}{r}\bigl(\delta_{il} - \hat{n}_i\,\hat{n}_l\bigr).
\label{eq:divG}
\end{equation}
Integrating:
\begin{equation}
  \int G_{ij,jl}\,d^3x
  = \delta_{il} \int_0^a r^2
    \left[\frac{4\pi}{3}\,q' + \frac{8\pi}{3}\,\frac{q}{r}\right] dr
  = (A + 4B)\,\delta_{il}.
\label{eq:contr2}
\end{equation}
Solving the $2\times 2$ system \eqref{eq:contr1}--\eqref{eq:contr2} yields
$A$ and~$B$, with only 1D radial integrals for SymPy to evaluate.


\section{Analytical Volume Integrals}

\subsection{Green's Tensor Integral}

Using the decomposition of Section~4, the volume integral of the Green's
tensor reduces to:
\begin{equation}
  \int_{\text{sphere}} G_{ij}\,d^3x
  = \delta_{ij} \int_0^a r^2
    \left[4\pi\,f(r) + \frac{4\pi}{3}\,g(r)\right] dr
  \equiv \Gamma_0\,\delta_{ij}.
\label{eq:Gamma0}
\end{equation}
After exact radial evaluation by SymPy:
\begin{equation}
  \Gamma_0 = \frac{%
    -3\alpha\beta
    + 2\alpha\bigl(\beta - i\omega a\bigr)\,e^{i\omega a/\beta}
    + \beta\bigl(\alpha - i\omega a\bigr)\,e^{i\omega a/\alpha}
  }{3\alpha\beta\,\omega^2\rho}.
\label{eq:Gamma0result}
\end{equation}

\subsection{Second-Derivative Integral}

As derived in Section~4.4, the volume integral of $G_{ij,kl}$ has the
isotropic structure
\begin{equation}
  \int_{\text{sphere}} G_{ij,kl}(\bx)\,d^3x
  = A\,\delta_{ij}\delta_{kl}
  + B\,(\delta_{ik}\delta_{jl} + \delta_{il}\delta_{jk}),
\label{eq:ABiso}
\end{equation}
with $A$ and $B$ determined by the two contractions
\eqref{eq:contr1}--\eqref{eq:contr2}.  Defining
$\xi_P = \omega a/\alpha$ and $\xi_S = \omega a/\beta$, the exact
closed-form results are:
\begin{equation}
  A = \frac{%
    2\alpha^3\beta\bigl(1-e^{i\xi_S}\bigr)
    + 4i\alpha^3\omega a\,e^{i\xi_S}
    + 3\alpha\beta^3\bigl(1-e^{i\xi_P}\bigr)
    + i\beta^3\omega a\,e^{i\xi_P}
  }{15\alpha^3\beta^3\rho},
\label{eq:Aresult}
\end{equation}
\begin{equation}
  B = \frac{%
    2\alpha^3\beta\bigl(1-e^{i\xi_S}\bigr)
    - i\alpha^3\omega a\,e^{i\xi_S}
    - 2\alpha\beta^3\bigl(1-e^{i\xi_P}\bigr)
    + i\beta^3\omega a\,e^{i\xi_P}
  }{15\alpha^3\beta^3\rho}.
\label{eq:Bresult}
\end{equation}

\smallskip
\noindent
\textit{Note:}  The parity argument shows $\int H_{ijk}\,d^3x = 0$, so the
stiffness contrast enters only through the strain equation.


\section{T-Matrix Coupling Coefficients}

The T-matrix tensor, defined by contracting the symmetrised
second-derivative integral with the stiffness contrast, is isotropic:
\begin{equation}
  T_{mnlp} = T_1\,\delta_{mn}\delta_{lp}
  + T_2\,(\delta_{ml}\delta_{np} + \delta_{mp}\delta_{nl}),
\label{eq:Tiso}
\end{equation}
where:
\begin{multline}
  T_1 = \frac{1}{15\alpha^3\beta^3\rho}\Bigl[
    \Delta\lambda\Bigl(
      10\alpha^3\beta\bigl(1-e^{i\xi_S}\bigr)
      + 5\alpha\beta^3\bigl(e^{i\xi_P}-1\bigr)
      + 5i\beta^3\omega a\,e^{i\xi_S}
    \Bigr) \\
    + \Delta\mu\Bigl(
      4\alpha^3\beta\bigl(1-e^{i\xi_S}\bigr)
      - 2i\alpha^3\omega a\,e^{i\xi_S}
      + 4\alpha\beta^3\bigl(e^{i\xi_P}-1\bigr)
      + 2i\beta^3\omega a\,e^{i\xi_P}
    \Bigr)
  \Bigr],
\label{eq:T1}
\end{multline}
with $\xi_P = \omega a/\alpha$ and $\xi_S = \omega a/\beta$.
\begin{equation}
  T_2 = \frac{\Delta\mu}{15\alpha^3\beta^3\rho}\Bigl[
    4\alpha^3\beta\bigl(1-e^{i\xi_S}\bigr)
    + 3i\alpha^3\omega a\,e^{i\xi_S}
    + \alpha\beta^3\bigl(1-e^{i\xi_P}\bigr)
    + 2i\beta^3\omega a\,e^{i\xi_P}
  \Bigr].
\label{eq:T2}
\end{equation}


\section{Self-Consistent Solution}

Under the assumption that the wavefield is approximately constant inside
the small sphere, the Lippmann--Schwinger equation reduces to a
self-consistent algebraic system that decouples into three independent
scalar equations for displacement, dilatation, and deviatoric strain, with
amplification factors:
\begin{equation}
  A_u = \frac{1}{1-\omega^2\Delta\rho\,\Gamma_0}, \qquad
  A_\theta = \frac{1}{1-3T_1-2T_2}, \qquad
  A_e = \frac{1}{1-2T_2}.
\label{eq:amplsph}
\end{equation}


\section{Effective (Renormalised) Contrasts --- The T-Matrix}

The scattered field takes the Born form with renormalised contrasts:
\begin{equation}
  u_i^{\text{scat}}(\bx)
  = V\Bigl[
    H_{ijk}(\bx,\bx_s)\,\bigl(
      \Delta\lambda^*\,\theta^0\,\delta_{jk}
      + 2\Delta\mu^*\,\varepsilon_{jk}^0\bigr)
    + \omega^2\,G_{ij}(\bx,\bx_s)\,\Delta\rho^*\,u_j^0
  \Bigr],
\label{eq:scatBorn}
\end{equation}
where $V = \tfrac{4}{3}\pi a^3$ is the sphere volume.  The effective
contrasts are:

\medskip
\noindent\textbf{Density:}
\begin{equation}
  \Delta\rho^* = \frac{\Delta\rho}{1-\omega^2\Delta\rho\,\Gamma_0}.
\label{eq:Drhostar}
\end{equation}

\noindent\textbf{Shear modulus:}
\begin{equation}
  \Delta\mu^* = \frac{\Delta\mu}{1-2T_2}.
\label{eq:Dmustar}
\end{equation}

\noindent\textbf{Lam\'e parameter:}
\begin{equation}
  \Delta\lambda^* = \frac{\Delta\lambda + \tfrac{2}{3}\Delta\mu}
  {1 - 3T_1 - 2T_2}
  - \frac{\tfrac{2}{3}\Delta\mu}{1 - 2T_2}.
\label{eq:Dlamstar}
\end{equation}


\section{Born Approximation Verification}

In the Born (point-scatterer) limit $a\to 0$, all amplification factors
tend to unity:
\begin{equation}
  \lim_{a\to 0}\Delta\rho^* = \Delta\rho, \qquad
  \lim_{a\to 0}\Delta\lambda^* = \Delta\lambda, \qquad
  \lim_{a\to 0}\Delta\mu^* = \Delta\mu.
\label{eq:Bornlim}
\end{equation}
All three limits have been verified analytically with SymPy.


\section{Numerical Example}

Background: $V_P = 5000$~m/s, $V_S = 3000$~m/s,
$\rho = 2500$~kg/m$^3$.  Contrasts: $\Delta\lambda = 2$~GPa,
$\Delta\mu = 1$~GPa, $\Delta\rho = 100$~kg/m$^3$.
Frequency $f = 10$~Hz, sphere radius $a = 10$~m.

\medskip
\begin{center}
\begin{tabular}{ll}
\toprule
\textbf{Quantity} & \textbf{Value} \\
\midrule
$k_P a$ & $0.1257$ \\
$k_S a$ & $0.2094$ \\
\midrule
$\Gamma_0$ & $1.730890\mathrm{e}{-09} + 2.282529\mathrm{e}{-10}\,i$ \\
$A$        & $-3.781591\mathrm{e}{-13} + 9.278517\mathrm{e}{-13}\,i$ \\
$B$        & $2.248959\mathrm{e}{-13} - 1.438708\mathrm{e}{-12}\,i$ \\
$T_1$      & $1.492641\mathrm{e}{-03} - 1.253138\mathrm{e}{-02}\,i$ \\
$T_2$      & $-1.532633\mathrm{e}{-04} - 5.108564\mathrm{e}{-04}\,i$ \\
$A_u$      & $1.000684\mathrm{e}{+00} + 9.023389\mathrm{e}{-05}\,i$ \\
$A_\theta$ & $1.002681\mathrm{e}{+00} - 3.888157\mathrm{e}{-02}\,i$ \\
$A_e$      & $9.996925\mathrm{e}{-01} - 1.021086\mathrm{e}{-03}\,i$ \\
$\Delta\rho^*$     & $1.000684\mathrm{e}{+02} + 9.023389\mathrm{e}{-03}\,i$ \\
$\Delta\lambda^*$  & $2.007355\mathrm{e}{+09} - 1.030035\mathrm{e}{+08}\,i$ \\
$\Delta\mu^*$      & $9.996925\mathrm{e}{+08} - 1.021086\mathrm{e}{+06}\,i$ \\
\bottomrule
\end{tabular}
\end{center}

% =============================================================================
%  Part II: Cubic Scatterer  (Sections 11--16, 18)
% =============================================================================

\part*{Part II: Cubic Scatterer}

\section{Motivation: Why Cubes?}

The T-matrix for a cube is fundamentally more useful for effective medium
theory than the T-matrix for a sphere, because cubes tile three-dimensional
space.

A sphere is the natural scatterer for studying isolated inclusions in a
matrix, but when the goal is to model a heterogeneous medium as a
homogenised effective material, one must partition the entire volume into
non-overlapping cells.  Spheres cannot fill space without gaps: even the
densest packing of equal spheres leaves roughly 26\% void \citep{Hales2005}.
Cubes, by contrast, are space-filling: they tesselate $\mathbb{R}^3$ with
zero overlap and zero void.

This simple geometric fact has far-reaching consequences for effective
medium theory (see Section~17).

\subsection{Semi-Analytical Approach}

For the sphere, the radial decomposition of $G_{ij}$ and the angular
integral identities allow exact closed-form evaluation of all volume
integrals.  For the cube, no such radial-angular factorisation exists.
Instead, we use a semi-analytical approach:

\begin{enumerate}
\item Taylor-expand the Green's tensor $G_{ij}(\bx)$ about the centre of
  the cube ($\bx=\mathbf{0}$) as a polynomial in $(x_1,x_2,x_3)$, keeping
  terms through order $(\omega a)^7$.

\item Integrate each monomial $x_1^{p_1} x_2^{p_2} x_3^{p_3}$ over the
  cube $[-a,a]^3$ analytically using the factorisation:
\begin{equation}
  \int_{-a}^{a}\!\int_{-a}^{a}\!\int_{-a}^{a}
  x_1^{p_1} x_2^{p_2} x_3^{p_3}\,d^3x
  = \prod_{m=1}^{3} \frac{2\,a^{p_m+1}}{p_m+1}
  \qquad(\text{all } p_m \text{ even; zero otherwise}).
\label{eq:cubemono}
\end{equation}
\end{enumerate}
This gives the same leading-order behaviour as the sphere (the $O(ka)^3$
Rayleigh terms), plus a cubic-anisotropy correction at $O(ka)^7$ that has
no analogue for the sphere.


\section{Cubic Anisotropy Tensor}

The critical difference between sphere and cube appears in the fourth
moment of the scatterer volume.  For any centrosymmetric shape:
\begin{equation}
  \int_D x_i\,x_j\,x_k\,x_l\,d^3x
  = P\,(\delta_{ij}\delta_{kl} + \delta_{ik}\delta_{jl}
    + \delta_{il}\delta_{jk})
  + Q\,E_{ijkl},
\label{eq:4thmoment}
\end{equation}
where $E_{ijkl} = \sum_m \delta_{im}\delta_{jm}\delta_{km}\delta_{lm}$ is
the \emph{cubic anisotropy tensor} (equals~1 when all four indices are
equal, zero otherwise).

\textbf{Sphere:}  By spherical symmetry, $Q=0$.  The fourth moment is
purely isotropic.

\textbf{Cube} $[-a,a]^3$:  Direct computation gives
\begin{equation}
  \int x_1^4\,d^3x = \frac{8a^7}{5}, \qquad
  \int x_1^2 x_2^2\,d^3x = \frac{8a^7}{9},
\label{eq:cube4th}
\end{equation}
so $P = 8a^7/9$ and $Q = -16a^7/15 \neq 0$.  The cube has genuine cubic
anisotropy.


\section{Volume Integrals (Cube)}

\subsection{Green's Tensor Integral $\Gamma_0^{\text{cube}}$}

The zeroth-order integral
$\int_{\text{cube}} G_{ij}\,d^3x = \Gamma_0^{\text{cube}}\,\delta_{ij}$
is still isotropic (only the second moment enters, and
$\int x_i x_j\,d^3x \propto \delta_{ij}$ for the cube).  The result
differs from the sphere only through the geometry of the integration domain.

\subsection{Second-Derivative Integral: Three Parameters}

For the cube, the second-derivative integral acquires a third independent
tensor structure:
\begin{equation}
  \int_{\text{cube}} G_{ij,kl}\,d^3x
  = A^c\,\delta_{ij}\delta_{kl}
  + B^c\,(\delta_{ik}\delta_{jl} + \delta_{il}\delta_{jk})
  + C^c\,E_{ijkl},
\label{eq:ABCcube}
\end{equation}
where $C^c$ is the cubic anisotropy coefficient.  Computing the three
independent components:
\begin{equation}
  A^c = I_{1122}, \qquad
  B^c = I_{1212}, \qquad
  C^c = I_{1111} - I_{1122} - 2\,I_{1212}.
\label{eq:ABCdef}
\end{equation}
SymPy computes $C^c = O\bigl((\omega a)^7\bigr)$, confirming it arises
from the fourth-moment anisotropy of the cube.


\section{T-Matrix Coupling Coefficients (Cube)}

The T-matrix tensor for the cube has cubic (not isotropic) symmetry:
\begin{equation}
  T_{mnlp}^c = T_1^c\,\delta_{mn}\delta_{lp}
  + T_2^c\,(\delta_{ml}\delta_{np} + \delta_{mp}\delta_{nl})
  + T_3^c\,E_{mnlp}.
\label{eq:Tcube}
\end{equation}
The third coefficient $T_3^c$ arises from $C^c$ and vanishes for a sphere
($C=0 \Rightarrow T_3=0$).


\section{Self-Consistent Solution: The $6\times 6$ Voigt Matrix
  and Its Eigenvalue Decomposition}

\subsection{The Voigt Representation}

The self-consistent strain equation $\varepsilon_{mn} = \varepsilon_{mn}^0
+ T_{mnlp}\,\varepsilon_{lp}$ is a $3\times 3$ symmetric tensor equation.
In Voigt notation, the six independent strain components
\begin{equation}
  \boldsymbol{\varepsilon}_V
  = [\varepsilon_{11},\;\varepsilon_{22},\;\varepsilon_{33},\;
     2\varepsilon_{23},\;2\varepsilon_{13},\;2\varepsilon_{12}]^T
\label{eq:Voigt}
\end{equation}
are mapped to a 6-vector, and $T_{mnlp}$ becomes a $6\times 6$ matrix
$\bT_V$.

For the cubic T-tensor~\eqref{eq:Tcube}, $\bT_V$ has the block-diagonal
structure:
\begin{equation}
  \bT_V = \begin{pmatrix} \bD & \mathbf{0} \\ \mathbf{0} & \bS \end{pmatrix},
\label{eq:TVblock}
\end{equation}
where the $3\times 3$ diagonal-strain block is
\begin{equation}
  \bD = T_1\,\bJ + (2T_2 + T_3)\,\bI_3
  = \begin{pmatrix}
    T_1+2T_2+T_3 & T_1 & T_1 \\
    T_1 & T_1+2T_2+T_3 & T_1 \\
    T_1 & T_1 & T_1+2T_2+T_3
  \end{pmatrix}
\label{eq:Dblock}
\end{equation}
(with $\bJ$ the $3\times 3$ matrix of ones) and the $3\times 3$
off-diagonal shear block is
\begin{equation}
  \bS = 2T_2\,\bI_3.
\label{eq:Sblock}
\end{equation}

The physical reason for the block-diagonal structure is that the cubic
anisotropy tensor $E_{mnlp}$ only contributes when $m=n$ and $l=p$ (all
indices equal).  When either pair is off-diagonal ($m\neq n$ or $l\neq p$),
$E_{mnlp}=0$, so the $T_3$ term drops out and the off-diagonal shear
components decouple with the isotropic eigenvalue $2T_2$.

\subsection{Eigenvalue Decomposition: 4 Scalar Amplification Factors}

The self-consistent equation $\boldsymbol{\varepsilon} =
\boldsymbol{\varepsilon}^0 + \bT_V\,\boldsymbol{\varepsilon}$ has the
formal solution
$\boldsymbol{\varepsilon} = (\bI_6 - \bT_V)^{-1}\,\boldsymbol{\varepsilon}^0$.
Instead of inverting the full $6\times 6$ matrix symbolically (which
involves extremely expensive \texttt{simplify()} calls on nested rational
expressions), we exploit the block structure to reduce everything to four
scalar equations.

\textbf{Block $\bS$ (off-diagonal shear):}  Already diagonal.  The
amplification factor is
\begin{equation}
  A_e^{(\text{off})} = \frac{1}{1-2T_2}.
\label{eq:Aeoff}
\end{equation}

\textbf{Block $\bD$ (diagonal strains):}  The matrix
$\bD = T_1\,\bJ + (2T_2+T_3)\,\bI_3$ has the well-known eigenstructure:

\begin{itemize}
\item \textbf{Dilatation mode:}  eigenvector $[1,1,1]^T$ (uniform
  expansion), eigenvalue $3T_1 + 2T_2 + T_3$:
\begin{equation}
  A_\theta = \frac{1}{1-3T_1-2T_2-T_3}.
\label{eq:Athcube}
\end{equation}

\item \textbf{Diagonal deviatoric modes:}  eigenvectors $[1,-1,0]^T$ and
  $[1,1,-2]^T$ (traceless diagonal strain), eigenvalue $2T_2+T_3$:
\begin{equation}
  A_e^{(\text{diag})} = \frac{1}{1-2T_2-T_3}.
\label{eq:Aediag}
\end{equation}
\end{itemize}
Plus the displacement amplification (unchanged from the sphere):
\begin{equation}
  A_u = \frac{1}{1-\omega^2\Delta\rho\,\Gamma_0^{\text{cube}}}.
\label{eq:Aucube}
\end{equation}

For the sphere ($T_3=0$): $A_e^{(\text{off})} = A_e^{(\text{diag})} \equiv
A_e$ and we recover the 3 amplification factors of Section~8.


\section{Effective Contrasts (Cube): Scalar Derivation}

The effective stiffness contrast is
$\Delta c^*_{mnlp} = \Delta c_{mnuv}\,A_{uv,lp}$, where $A$ is the
amplification operator.  We avoid the full $6\times 6$ matrix product by
decomposing the incident strain into the three eigenspaces of~$A$ and
applying each scalar amplification factor separately.

The isotropic contrast $\Delta c_{mnuv}\,\varepsilon_{uv}^*
= \Delta\lambda\,\theta^*\,\delta_{mn} + 2\Delta\mu\,\varepsilon_{mn}^*$.
Substituting the amplified strain:
\begin{equation}
  \varepsilon_{mn}^*
  = \underbrace{\frac{A_\theta\,\theta^0}{3}\,\delta_{mn}}_{\text{dilatation}}
  + \underbrace{A_e^{(\text{diag})}\,e_{mn}^{0,\text{diag}}}_{\text{diag.\ deviatoric}}
  + \underbrace{A_e^{(\text{off})}\,e_{mn}^{0,\text{off}}}_{\text{off-diag.\ shear}}
\label{eq:ampstrain}
\end{equation}
gives, after algebra:

\medskip
\noindent\textbf{Effective density:}
\begin{equation}
  \Delta\rho^{*,c} = \frac{\Delta\rho}
  {1-\omega^2\Delta\rho\,\Gamma_0^{\text{cube}}}.
\label{eq:Drhostar_c}
\end{equation}

\noindent\textbf{Off-diagonal effective shear modulus:}
\begin{equation}
  \Delta\mu^{*,\text{off}} = \frac{\Delta\mu}{1-2T_2^c}.
\label{eq:Dmuoff}
\end{equation}

\noindent\textbf{Diagonal effective shear modulus:}
\begin{equation}
  \Delta\mu^{*,\text{diag}} = \frac{\Delta\mu}{1-2T_2^c-T_3^c}.
\label{eq:Dmudiag}
\end{equation}

\noindent\textbf{Effective Lam\'e parameter:}
\begin{equation}
  \Delta\lambda^{*,c}
  = \frac{\Delta\lambda + \tfrac{2}{3}\Delta\mu}
         {1-3T_1^c-2T_2^c-T_3^c}
  - \frac{\Delta\mu}{3}
    \left(\frac{1}{1-2T_2^c-T_3^c} + \frac{1}{1-2T_2^c}\right).
\label{eq:Dlamstar_c}
\end{equation}

The cubic anisotropy of the effective shear modulus is:
\begin{equation}
  \Delta\mu^{*,\text{diag}} - \Delta\mu^{*,\text{off}}
  = \frac{\Delta\mu\,T_3^c}
  {(1-2T_2^c)(1-2T_2^c-T_3^c)}
  = O\bigl((\omega a)^7\bigr).
\label{eq:cubicaniso}
\end{equation}

% =============================================================================
%  Part III: From Single Scatterer to Effective Medium  (Sections 17--18)
% =============================================================================

\part*{Part III: From Single Scatterer to Effective Medium}

\section{Homogenisation and Effective Medium Theory}

\subsection{The Space-Filling Requirement}

The ultimate goal of computing single-scatterer T-matrices is to build
effective medium theories for heterogeneous elastic materials.  The idea is
simple: model the heterogeneous medium as a collection of scatterers, each
characterised by its T-matrix, and then average over the ensemble to obtain
effective material properties.

The critical requirement for such a homogenisation procedure is that the
scatterers must \emph{tile} the medium completely.  If they do not, the
material between scatterers is unmodelled, leading to systematic
errors \citep{Willis1977,Hashin1983}.

Cubes are the simplest three-dimensional shape that tiles $\mathbb{R}^3$
without gaps or overlaps.  This was famously proven to be optimal for convex
shapes by the Keller--Minkowski theorem, and the practical importance for
effective medium theory was recognised early on by Hashin and
Shtrikman \citep{HashinShtrikman1962a,HashinShtrikman1962b,HashinShtrikman1963},
who used cubic partitions as the foundation of their variational bounds on
effective elastic moduli.

\subsection{From T-Matrix to Effective Moduli}

Given the cubic T-matrix~\eqref{eq:Tcube}, we can construct the effective
elastic moduli of a composite material as follows.  Consider a background
medium partitioned into $N$ non-overlapping cubes, each with (possibly
different) material properties.  In the long-wavelength limit, the
effective stiffness tensor satisfies \citep{ZellerDederichs1973,Willis1977}:
\begin{equation}
  C_{ijkl}^{\text{eff}} = C_{ijkl}^{(0)}
  + \frac{1}{V}\sum_{n=1}^{N}
    \Delta c_{ijuv}^{(n)}\,A_{uv,kl}^{(n)},
\label{eq:Ceff}
\end{equation}
where $C^{(0)}$ is the background stiffness, $\Delta c^{(n)}$ is the
contrast of the $n$-th cube, and $A^{(n)}$ is the amplification operator
derived from its T-matrix.  In the dilute limit (non-interacting
scatterers), each $A^{(n)}$ is computed independently as in Section~15.

\subsection{Self-Consistent Effective Medium Schemes}

More sophisticated approximations account for scatterer-scatterer
interactions:

\textbf{Coherent Potential Approximation (CPA):}  Each cube is embedded in
the (unknown) effective medium rather than the background, leading to a
self-consistent equation for $C^{\text{eff}}$ \citep{ZellerDederichs1973,Gubernatis1977}.
The T-matrix formulation is particularly natural here: the CPA condition
is simply that the average T-matrix vanishes, $\langle T\rangle = 0$.

\textbf{Differential Effective Medium (DEM):}  Scatterers are added
incrementally to the background, with the effective medium updated at each
step \citep{Norris1985}.  The cubic T-matrix provides the ``update kernel''
for each infinitesimal addition.

\textbf{Kramers--Kronig consistency:}  Because the T-matrix is derived
from the full self-consistent Lippmann--Schwinger equation (not the Born
approximation), the resulting effective moduli are causal and satisfy
Kramers--Kronig relations \citep{Willis1981}, making them suitable for
broadband modelling.

\subsection{Cubic Symmetry of the Effective Medium}

Even when the individual cubes have isotropic material properties
($\Delta\lambda$, $\Delta\mu$, $\Delta\rho$), the effective medium inherits
\emph{cubic symmetry} from the scatterer geometry.  This is a direct
consequence of $T_3\neq 0$ in~\eqref{eq:Tcube}: the effective shear
modulus splits into $\mu^{*,\text{diag}}$ and $\mu^{*,\text{off}}$, giving
the effective stiffness tensor three independent elastic constants (bulk
modulus, and two shear moduli) rather than the two of an isotropic medium.

The cubic anisotropy parameter~\eqref{eq:cubicaniso} quantifies this
splitting.  At low frequencies it is $O(\omega a)^7$---very small---but it
grows rapidly with frequency, becoming significant as $ka$ approaches unity.

\subsection{Connection to Hashin--Shtrikman Bounds}

The classical Hashin--Shtrikman variational
bounds \citep{HashinShtrikman1962a,HashinShtrikman1962b,HashinShtrikman1963}
on effective elastic moduli are derived using a comparison medium and a
polarisation stress field.  The bounds are achieved by specific
microstructural geometries: the Hashin assemblage (concentric spheres)
achieves the bulk modulus bound, but the shear modulus bound requires more
complex geometries.

The cubic T-matrix approach provides a complementary, wave-based route to
effective properties that naturally includes frequency dependence and
finite-wavelength corrections absent from static variational bounds.  In
the zero-frequency limit, the cubic effective moduli computed here should
lie within the Hashin--Shtrikman bounds, providing a consistency check.


\section{Born Verification (Cube)}

In the Born limit $a\to 0$:
\begin{equation}
  \lim_{a\to 0}\Delta\rho^{*,c} = \Delta\rho, \qquad
  \lim_{a\to 0}\Delta\lambda^{*,c} = \Delta\lambda, \qquad
  \lim_{a\to 0}\Delta\mu^{*,\text{diag}} = \Delta\mu, \qquad
  \lim_{a\to 0}\Delta\mu^{*,\text{off}} = \Delta\mu.
\label{eq:Bornlim_cube}
\end{equation}
The cubic anisotropy vanishes:
$\lim_{a\to 0}(\Delta\mu^{*,\text{diag}} - \Delta\mu^{*,\text{off}}) = 0$.
All limits verified analytically with SymPy.

% =============================================================================
%  Part IV: The Global T-Matrix  (Sections 19--22)
% =============================================================================

\part*{Part IV: The Global T-Matrix}

\section{From Local Site T-Matrix to Global Multiple Scattering}

The T-matrix $\bT_0$ derived in Parts~I and~II is a \emph{local site} (or
single-scatterer) quantity: it describes the self-consistent response of one
cubic cell in isolation, embedded in a homogeneous reference medium.  To
model a heterogeneous medium composed of many such cells, we must account
for the multiple scattering between sites---each cube scatters incident
waves, which propagate to neighbouring cubes, scatter again, and so on.

The standard framework for this is the Foldy--Lax multiple-scattering
theory \citep{Foldy1945,Lax1952}, recast in the T-matrix formalism by
\citet{Waterman1976} and applied to elastic media by Zeller \&
\citet{ZellerDederichs1973},
\citet{Gubernatis1977}, and \citet{Jakobsen2012}.

\subsection{The Lippmann--Schwinger Equation on a Lattice}

Partition the medium into $N$ non-overlapping cubes centred at positions
$\{\bx_n\}_{n=1}^N$, each with its own material contrasts
$\Delta c^{(n)}$, $\Delta\rho^{(n)}$.  The total scattered field at site~$m$
due to all other sites is:
\begin{equation}
  u_i^{\text{tot}}(\bx_m)
  = u_i^0(\bx_m)
  + \sum_{n\neq m}
    G_{ij}^{(0)}(\bx_m - \bx_n)\;\bT_0^{(n)}\cdot u_j^{\text{tot}}(\bx_n),
\label{eq:latticeLS}
\end{equation}
where $G_{ij}^{(0)}(\bx_m - \bx_n)$ is the background (reference medium)
Green's tensor evaluated at the vector separation between the centres of
cubes~$m$ and~$n$, and $\bT_0^{(n)}$ is the local site T-matrix for the
$n$-th cube (computed as in Part~II).

This is the elastic analogue of the Foldy--Lax
equations \citep{Foldy1945,Lax1952}.  In compact operator notation, writing
$\bu = [u(\bx_1),\ldots,u(\bx_N)]^T$ and assembling the inter-site
Green's tensors into a matrix~$\bG_0$:
\begin{equation}
  \bu = \bu^0 + \bG_0\,\bT_0\,\bu,
\label{eq:compactLS}
\end{equation}
where $\bT_0 = \mathrm{diag}(\bT_0^{(1)},\ldots,\bT_0^{(N)})$ is the
block-diagonal matrix of local site T-matrices.

\subsection{The Global T-Matrix}

Solving~\eqref{eq:compactLS} for~$\bu$ gives
$\bu = (\bI - \bG_0\,\bT_0)^{-1}\,\bu^0$.  The global T-matrix is
defined as the operator that maps the incident field to the total scattered
field:
\begin{equation}
  \bT = \bT_0\,(\bI - \bG_0\,\bT_0)^{-1}.
\label{eq:Tglobal}
\end{equation}
Equivalently, by pre-multiplying with $\bG_0$:
\begin{equation}
  \bG_0\,\bT
  = \bG_0\,\bT_0\,(\bI - \bG_0\,\bT_0)^{-1}
  = (\bI - \bG_0\,\bT_0)^{-1} - \bI.
\label{eq:G0T}
\end{equation}

\subsection{Physical Interpretation: The Neumann Series}

When $\|\bG_0\,\bT_0\| < 1$ (weak scattering), the inverse can be
expanded as a Neumann series:
\begin{equation}
  \bT = \bT_0 + \bT_0\,\bG_0\,\bT_0
  + \bT_0\,\bG_0\,\bT_0\,\bG_0\,\bT_0 + \cdots
\label{eq:Neumann}
\end{equation}

Each term has a transparent physical interpretation (Figure~1):

\begin{itemize}
\item \textbf{Zeroth order} $\bT_0$: Single scattering at each site (the
  Born--T-matrix approximation).

\item \textbf{First order} $\bT_0\,\bG_0\,\bT_0$: A wave scatters at
  site~$n$ ($\bT_0^{(n)}$), propagates to site~$m$ ($\bG_0$), and scatters
  again ($\bT_0^{(m)}$).

\item \textbf{$k$-th order:}  Chains of $(k+1)$ scattering events
  connected by $k$ propagation steps.
\end{itemize}
The full inverse $(\bI - \bG_0\,\bT_0)^{-1}$ sums this series to all
orders, capturing the complete multiple-scattering response.

\begin{quote}
\textit{[Schematic to be added.]}  A pictorial representation showing:
(a)~a regular lattice of cubes, each labelled with its local T-matrix
$\bT_0^{(n)}$;
(b)~arrows representing the inter-site Green's tensor
$G_{ij}^{(0)}(\bx_m-\bx_n)$ connecting cube centres;
(c)~the multiple-scattering chain: incident wave $\to$ scatter at site~$n$
$\to$ propagate $\to$ scatter at site~$m$ $\to$ propagate $\to\cdots$
\end{quote}

\noindent\textbf{Figure~1:}  Schematic of the global T-matrix formalism.
Left: the medium is partitioned into cubes, each with a local T-matrix
$\bT_0^{(n)}$.  Right: the inter-site Green's tensor $\bG_0$ propagates the
scattered field between cube centres, and the global T-matrix~\eqref{eq:Tglobal}
sums all multiple-scattering paths to all orders.


\section{The Inter-Site Green's Tensor $\bG_0$}

\subsection{Homogeneous Reference Medium}

In a homogeneous isotropic reference medium, the inter-site Green's tensor
is the free-space Green's tensor evaluated at the separation vector:
\begin{equation}
  G_{ij}^{(0)}(\bx_m - \bx_n)
  = f(|\bx_m-\bx_n|)\,\delta_{ij}
  + g(|\bx_m-\bx_n|)\,\hat{r}_{mn,i}\,\hat{r}_{mn,j},
\label{eq:G0homog}
\end{equation}
where $f(r)$ and $g(r)$ are the radial functions from Section~3, and
$\hat{r}_{mn,i} = (x_{m,i}-x_{n,i})/|\bx_m-\bx_n|$ is the unit vector
from cube~$n$ to cube~$m$.

For a regular cubic lattice with spacing $2a$ (so that adjacent cube faces
are touching), the inter-site vectors are
$\bx_m - \bx_n = 2a\,(p,q,r)$ for integers $(p,q,r)$ with
$(p,q,r)\neq(0,0,0)$.

\subsection{Layered Reference Medium: The Kennett Propagator}

For geophysical applications, the background medium is not homogeneous but
\emph{stratified}---consisting of horizontal layers with depth-dependent
properties.  In this case, the inter-site Green's tensor must be replaced
by the Green's tensor for the layered reference medium.

\citet{Kennett1983} showed that the elastic Green's tensor for a
plane-layered medium can be computed efficiently using the reflectivity
method: the wavefield is decomposed into plane waves (via horizontal
wavenumber integration), each of which propagates through the layer stack
via a system of reflection and transmission matrices.  The key recursive
relations are:

\begin{enumerate}
\item \textbf{Layer propagator:}  For a homogeneous layer of thickness~$h$
  with P- and S-wave vertical slownesses $q_\alpha$ and~$q_\beta$, the
  phase propagation across the layer is given by
  $\mathbf{E}(h) = \mathrm{diag}(e^{i\omega q_\alpha h},
  e^{i\omega q_\beta h})$.

\item \textbf{Interface coefficients:}  At each interface, Zoeppritz-type
  reflection and transmission coefficients $r_d, r_u, t_d, t_u$ (for both
  P--P, P--SV, SV--P, SV--SV conversions) relate the up- and down-going
  wave amplitudes.

\item \textbf{Kennett recursive algorithm:}  The generalised reflection
  matrix $R_D$ for a stack of layers is built recursively from the bottom
  up:
\begin{equation}
  R_D^{(j)} = r_d^{(j)}
  + t_u^{(j)}\,R_D^{(j+1)}\,
    \bigl(\bI - r_u^{(j)}\,R_D^{(j+1)}\bigr)^{-1}\,t_d^{(j)},
\label{eq:Kennett}
\end{equation}
where $r_d^{(j)}, t_d^{(j)}$ are the reflection/transmission matrices at
the $j$-th interface, and $R_D^{(j+1)}$ is the already-computed reflection
matrix for all layers below interface~$j$ \citep{Kennett1983,Kennett1979}.

\item \textbf{Green's tensor assembly:}  The frequency--wavenumber-domain
  Green's tensor is assembled from the reflection matrices and layer
  propagators.  A final inverse Hankel (or Fourier) transform over
  horizontal wavenumber yields
  $G_{ij}^{(0)}(\bx_m-\bx_n;\omega)$ in the space-frequency domain.
\end{enumerate}

\subsection{Significance for Heterogeneous Media Modelling}

The combination of the Kennett propagator for $\bG_0$ with the cubic local
T-matrix $\bT_0$ provides a powerful framework for modelling wave
propagation in realistic heterogeneous Earth models:

\begin{enumerate}
\item \textbf{Background model:}  A plane-layered velocity-density profile
  (e.g.\ from well logs or a 1D Earth model such as PREM or AK135) defines
  the reference medium.  The Green's tensor $\bG_0$ for this layered
  background is computed using the Kennett reflectivity method.

\item \textbf{3D heterogeneity:}  The departures from the layered background
  are partitioned into cubic cells, each with constant (but arbitrary)
  material contrasts $\Delta\lambda^{(n)}$, $\Delta\mu^{(n)}$,
  $\Delta\rho^{(n)}$.  The local T-matrix $\bT_0^{(n)}$ for each cube is
  computed analytically (Part~II)---it depends only on the cube's material
  contrasts and the local background properties.

\item \textbf{Multiple scattering:}  The global
  T-matrix~\eqref{eq:Tglobal} is assembled from the $\bT_0^{(n)}$ and
  $\bG_0$, and the total scattered field is obtained by applying $\bT$ to
  the incident field.
\end{enumerate}
This is essentially the approach of \citet{Jakobsen2012} and
\citet{JakobsenWu2016}, but with the important refinement
that the local T-matrix is computed self-consistently (not in the Born
approximation), allowing strong contrasts to be modelled accurately.


\section{Operator Structure and Symmetries}

\subsection{Block Structure of the Global System}

The global system~\eqref{eq:compactLS} has a natural block structure.  For
$N$ cubes, each described by 9 degrees of freedom (3 displacement + 6
strain components), the system matrices are of size $9N\times 9N$.  However,
the local T-matrix $\bT_0$ is block-diagonal (each block is $9\times 9$),
and the inter-site Green's tensor $\bG_0$ has translational structure (for a
regular lattice, $G_{mn}$ depends only on $\bx_m - \bx_n$).

For a regular cubic lattice, this translational invariance can be exploited:
the system~\eqref{eq:compactLS} becomes block-circulant and can be
diagonalised by a spatial Fourier transform over the lattice.  In Fourier
space, the global T-matrix at wavevector~$\mathbf{k}$ reduces to:
\begin{equation}
  \hat{\bT}(\mathbf{k})
  = \hat{\bT}_0\,\bigl(\bI - \hat{\bG}_0(\mathbf{k})\,\hat{\bT}_0\bigr)^{-1},
\label{eq:Tfourier}
\end{equation}
where $\hat{\bG}_0(\mathbf{k})
= \sum_{\bx_n\neq\mathbf{0}} \bG_0(\bx_n)\,e^{-i\mathbf{k}\cdot\bx_n}$
is the lattice Green's function (also called the structure factor of the
Green's tensor).

This Fourier decomposition reduces the $9N\times 9N$ system to $N$
independent $9\times 9$ systems, one for each wavevector in the first
Brillouin zone of the cubic lattice---a dramatic computational saving.

\subsection{Relationship to the Dyson Equation}

Equation~\eqref{eq:Tglobal} is the lattice analogue of the Dyson equation
from many-body theory.  Defining the self-energy (or mass operator) as
$\Sigma = \bT_0$, the effective Green's function of the heterogeneous
medium satisfies:
\begin{equation}
  \bG_{\text{eff}} = \bG_0 + \bG_0\,\Sigma\,\bG_{\text{eff}}
  = (\bI - \bG_0\,\Sigma)^{-1}\,\bG_0.
\label{eq:Dyson}
\end{equation}
The poles of $\bG_{\text{eff}}$ in the frequency--wavenumber domain give the
dispersion relations of the effective medium, including attenuation due to
scattering.  This connection to the Dyson equation was exploited by Zeller
\& \citet{ZellerDederichs1973} for elastic media and by
\citet{Jakobsen2012} for seismic modelling.


\section{Computational Methods for Inverting the Global System}

The central computational problem is to solve the linear system
\begin{equation}
  (\bI - \bG_0\,\bT_0)\,\bu = \bu^0,
\label{eq:globalsys}
\end{equation}
which is of dimension $9N\times 9N$ (three displacement and six independent
strain components per cube, $N$~cubes).  For realistic 3D models, $N$ may
range from $10^3$ (small exploration-scale problems) to $10^6$ or beyond
(crustal-scale), making the choice of solution method critical.

We survey five approaches, ranging from the simplest (direct inversion, the
Neumann series) to the most sophisticated (fast multipole methods).  Each
has a distinct domain of applicability.

\subsection{Direct Inversion}

The most straightforward approach is to form and factorise
$(\bI - \bG_0\,\bT_0)$ explicitly, using LU decomposition or a similar
dense solver.  This costs $O(N^3)$ in time and $O(N^2)$ in storage,
restricting it to small problems ($N\lesssim 10^3$).  However, it yields
the complete resolvent, from which the Green's function for all
source--receiver pairs can be extracted.  For benchmark problems and
validation of approximate methods, direct inversion remains indispensable.

\subsection{The Neumann (Born) Series}

The Neumann series~\eqref{eq:Neumann} is a fixed-point iteration:
\begin{equation}
  \bu^{(k+1)} = \bu^0 + \bG_0\,\bT_0\,\bu^{(k)}, \qquad
  \bu^{(0)} = \bu^0.
\label{eq:NeumannIter}
\end{equation}
It converges if and only if the spectral radius
$\rho(\bG_0\,\bT_0) < 1$.  In scattering theory this is the Born series,
and it is well known to diverge for strong material
contrasts \citep{Gubernatis1977}.  Even when it converges, the rate is
geometric with ratio~$\rho$, so many iterations may be needed for moderate
contrasts.  The Neumann series is therefore useful mainly as a conceptual
tool and as the starting point for more sophisticated iterative methods.

\subsection{Iterative Krylov Subspace Methods}

Krylov subspace methods---GMRES \citep{SaadSchultz1986},
BiCGSTAB \citep{vanderVorst1992}, and their variants---are the practical
workhorse for solving~\eqref{eq:globalsys}.  They require only the
matrix--vector product $\bv \mapsto \bG_0\,\bT_0\,\bv$, never the explicit
matrix, and achieve convergence in far fewer iterations than the Neumann
series because they minimise the residual over the full Krylov subspace
rather than iterating a single fixed-point map.

\textbf{GMRES} (Generalised Minimal Residual) \citep{SaadSchultz1986} is
the method of choice for non-symmetric systems.  At iteration~$k$ it
minimises $\|\mathbf{r}^{(k)}\|$ over the Krylov subspace
$\mathcal{K}_k = \mathrm{span}\{\mathbf{r}^{(0)},
\mathbf{A}\mathbf{r}^{(0)}, \ldots,
\mathbf{A}^{k-1}\mathbf{r}^{(0)}\}$, where
$\mathbf{A} = \bI - \bG_0\,\bT_0$.  Storage grows linearly with~$k$
(one vector per iteration), which can be controlled by restarting
(GMRES($m$)).

\textbf{BiCGSTAB} \citep{vanderVorst1992} is an alternative that avoids the
growing storage of GMRES at the cost of a less smooth convergence history.
Each iteration requires two matrix--vector products but only a fixed amount
of storage.

The matrix--vector product $\bG_0\,\bT_0\,\bv$ has a transparent physical
interpretation: (i)~\emph{scatter}: apply the block-diagonal $\bT_0$
to~$\bv$ (cost $O(N)$, since each $9\times 9$ block acts independently);
(ii)~\emph{propagate}: apply the dense inter-site Green's tensor $\bG_0$
(cost $O(N^2)$ naively, but reducible to $O(N\log N)$ with the methods of
Section~22.5).

\textbf{Preconditioning.}  The convergence rate of Krylov methods depends
on the spectral properties of $\bI - \bG_0\,\bT_0$.  Because the local
T-matrix $\bT_0$ already absorbs the strong single-site scattering
self-consistently, the remaining inter-site coupling $\bG_0\,\bT_0$ is
typically much weaker than the original material contrast, leading to
favourable convergence even without explicit preconditioning.  When
additional acceleration is needed, block-diagonal or block-Jacobi
preconditioners (using the diagonal blocks of $\bG_0$) are natural choices.

\textbf{Historical note.}  Krylov methods were applied to elastic-wave
scattering in layered media with superimposed spherical heterogeneities by
\citet{Nestor1996}, who solved the multiple-scattering series for the
global T-matrix using iterative methods with the Kennett reflectivity
method (Section~20.2) providing the inter-site Green's tensor $\bG_0$ and
spherical scatterer T-matrices for~$\bT_0$.  This combination demonstrated
the feasibility of modelling 3D heterogeneous media as perturbations to a
stratified background, but was limited by the isotropy of the spherical
T-matrix and the inability of spheres to tile space.  The cubic T-matrix
developed in Part~II removes both limitations, completing the bridge from
single-scatterer physics to rigorous homogenisation that was implicit in
the original formulation.

\subsection{The Renormalised Scattering Series}

\citet{JakobsenWu2016} proposed a renormalised scattering
series that overcomes the convergence limitations of the Born/Neumann
series for strong contrasts.  The idea is to split the total perturbation
into a ``strong'' part (absorbed into a modified reference medium or into
the local T-matrix) and a ``weak'' residual, then iterate on the residual
only.

Concretely, define a modified operator
$\tilde{\bG}_0 = \bG_0 - \bG_0^{\text{diag}}$, where
$\bG_0^{\text{diag}}$ contains only the self-interaction terms (which are
already accounted for in $\bT_0$).  The renormalised series iterates:
\begin{equation}
  \bu^{(k+1)} = \bu^0 + \tilde{\bG}_0\,\bT_0\,\bu^{(k)}.
\label{eq:renorm}
\end{equation}
Because $\tilde{\bG}_0$ excludes the (already-resummed) on-site terms, the
spectral radius $\rho(\tilde{\bG}_0\,\bT_0)$ is typically much smaller
than $\rho(\bG_0\,\bT_0)$, extending the domain of convergence to stronger
contrasts.  When combined with the self-consistent (non-Born) local
T-matrix, this approach effectively sums the dominant scattering
contributions exactly and treats only the residual inter-site coupling
perturbatively.

The renormalised series can also be viewed as a preconditioned Neumann
iteration, and it can be further accelerated by embedding it within a
Krylov solver.

\subsection{Fast Multipole and Hierarchical Matrix Methods}

The dominant cost of each Krylov iteration is the dense matrix--vector
product $\bG_0\,\bv$, which costs $O(N^2)$ when computed directly.  For
large~$N$, this can be reduced to $O(N\log N)$ or even $O(N)$ by
exploiting the spatial structure of the Green's tensor.

\textbf{Fast Multipole Method (FMM):}  The elastic Green's tensor
$G_{ij}^{(0)}(\bx)$ decays as $1/|\bx|$ and is smooth (low-rank) when
source and receiver are well separated.  The
FMM \citep{Rokhlin1985,GreengardRokhlin1987} exploits this by grouping
cubes into a hierarchical tree structure, computing far-field interactions
via multipole expansions, and only evaluating near-field interactions
directly.  For elastodynamics, the FMM was developed by
\citet{Fujiwara2000} and \citet{Chaillat2008},
achieving $O(N\log N)$ or $O(N)$ complexity per matrix--vector product.

\textbf{Hierarchical matrices ($\Hmat$-matrices):}  An alternative to the
FMM, $\Hmat$-matrices \citep{Hackbusch1999,Boerm2003} represent the
off-diagonal blocks of $\bG_0$ in low-rank factored form, again achieving
$O(N\log N)$ storage and matrix--vector products.  The advantage over FMM
is that $\Hmat$-matrices can also be inverted or LU-factored in
$O(N\log^2 N)$ time, providing a direct (non-iterative) solver.

\textbf{Hybrid FFT for layered $\bG_0$:}  When the reference medium is
layered (Section~20.2) and the cubes are arranged on a regular horizontal
grid (but at variable depths), the horizontal convolution structure of
$\bG_0$ can be exploited via 2D FFT.  Only the vertical coupling between
different depth levels requires explicit summation.  This hybrid approach
naturally suits the ``layered background + 3D perturbations'' geometry of
geophysical applications.


\subsection{Summary and Recommendations}

Table~1 summarises the computational methods and their domains of
applicability.

For the geophysical application envisioned here---3D heterogeneities
superimposed on a layered background via the Kennett propagator---the
recommended approach is a Krylov solver (GMRES or BiCGSTAB) with the
matrix--vector product $\bG_0\,\bT_0\,\bv$ accelerated by the hybrid
FFT/direct scheme described in Section~22.5.  The self-consistent local
T-matrix $\bT_0$ from Part~II serves as a natural preconditioner by
absorbing the strong single-site scattering, leaving only the weaker
inter-site coupling to be resolved iteratively.

\begin{table}[ht]
\centering
\begin{tabular}{llll}
\toprule
\textbf{Method} & \textbf{Cost} & \textbf{Requires} & \textbf{Best for} \\
\midrule
Direct inversion      & $O(N^3)$                       & Dense matrix              & Small $N$, benchmarks \\
Neumann series        & $O(kN^2)$                      & $\rho < 1$                & Weak scattering only \\
GMRES/BiCGSTAB        & $O(kN^2)$ or $O(kN\!\log\!N)$  & Mat-vec product           & General heterogeneity \\
Renormalised series   & $O(kN^2)$ or $O(kN\!\log\!N)$  & Self-consist.\ $T_0$      & Strong contrasts \\
FMM / $\Hmat$-matrices & $O(N\!\log\!N)$ per iter. & Spatial decay of $G_0$ & Large $N$, 3D \\
\bottomrule
\end{tabular}
\caption{Comparison of computational methods for solving the global T-matrix system~\eqref{eq:globalsys}.
Here $k$ denotes the number of iterations required for convergence.}
\end{table}

\bigskip
\noindent
\textit{Generated from \texttt{tmatrix\_analytic.py} and \texttt{tmatrix\_cube.py}
using SymPy symbolic computation.  All sphere integrals evaluated in exact closed
form; cube integrals evaluated semi-analytically via Taylor expansion to
$O(\omega a)^8$.}

% =============================================================================
%  Part V: Numerical Implementation  (Sections 23--26)
% =============================================================================

\part*{Part V: Numerical Implementation}

\section{Hybrid Computation Strategy}

The volume integrals derived in Parts~I and~II can be evaluated analytically
by SymPy (symbolic computer algebra) for both the sphere and cube geometries.
However, for deployment in a numerical wavefield-modelling code, a pure-NumPy
implementation is desirable.  This section describes the hybrid numerical
strategy used in the \texttt{CubicTMatrix} Python package.

\subsection{The Singularity Structure of the Integrands}

The two volume integrals that define the local site T-matrix have
fundamentally different singularity structures at the origin:

\textbf{Green's tensor integral} (Eq~\ref{eq:Gamma0} / Section~13.1): The integrand
$G_{ij}(\bx) \sim 1/r$ as $r \to 0$. In three dimensions,
$d^3x \sim r^2\,dr$, so the integral converges absolutely. Standard 3D
Gauss--Legendre quadrature evaluates it accurately.

\textbf{Second-derivative integral} (Eq~\ref{eq:ABCcube}): The integrand
$G_{ij,kl}(\bx) \sim 1/r^3$ as $r \to 0$. This is the static
Eshelby singularity. The volume integral \emph{diverges}---it must be
understood as a Cauchy principal value plus a shape-dependent depolarisation
tensor. Direct numerical quadrature captures the divergent static part with
an $n$-dependent finite remainder, producing incorrect results.

\subsection{The Hybrid Approach}

The solution is to use a different method for each integral:
\begin{itemize}
\item \textbf{$\Gamma_0$: full numerical quadrature.}  Vectorised 3D
  Gauss--Legendre quadrature over $[-a,a]^3$ with $n_{\mathrm{GL}} = 32$
  points per dimension.  The GL nodes never include the origin (they are
  the zeros of Legendre polynomials on $(-1,1)$), and the $1/r$
  singularity is integrable, so the integral converges rapidly.  This
  correctly captures both the real (static) and imaginary
  (frequency-dependent) contributions.

\item \textbf{$A^c$, $B^c$, $C^c$: polynomial Taylor expansion.}  We
  decompose the Green's tensor into a singular static part ($1/r^3$)
  and a smooth polynomial remainder.  Only the smooth remainder is
  integrated---the static singularity drops out automatically because it
  is excluded from the Taylor series of the frequency-dependent terms.
  The polynomial integrand is evaluated by GL quadrature, which
  integrates it exactly (to Taylor truncation order).
\end{itemize}
This hybrid strategy achieves the accuracy of the SymPy symbolic
computation (matching $T_1$ and $T_2$ to $\sim 10^{-8}$, $T_3$ to
$\sim 0.3\%$) while running in pure NumPy at negligible computational
cost.


%=====================================================================
\section{Smooth Green's Tensor Decomposition}
%=====================================================================

\subsection{Splitting off the Singularity}

The Green's tensor (Eq~\ref{eq:Grad}) can be decomposed as:
\begin{equation}
  G_{ij}(\bx)
  = G_{ij}^{\mathrm{static}}(\bx)
  + G_{ij}^{\mathrm{s}}(\bx),
\end{equation}
where the static part $G_{ij}^{\mathrm{static}} \sim 1/r$ produces the
divergent $1/r^3$ second derivatives, and the smooth part
$G_{ij}^{\mathrm{s}}$ is an entire function of $\bx$ (a convergent
power series).

The smooth part arises from the odd-power terms in the Taylor expansion of
$f(r)\cdot r$ and $g(r)\cdot r$ about $r=0$.  It has the structure:
\begin{equation}
  \boxed{%
    G_{ij}^{\mathrm{s}}(\bx)
    = \delta_{ij}\,\Phi(r^2) + x_i\,x_j\,\Psi(r^2)
  }
\end{equation}
where $\Phi(u) = \sum_n \phi_n\,u^n$ and $\Psi(u) = \sum_n \psi_n\,u^n$
are power series in $u = r^2$.

\subsection{Taylor Coefficients}

The Taylor coefficients are derived by expanding the exponentials
$e^{i\omega r/\alpha}$ and $e^{i\omega r/\beta}$ in the definitions of
$f(r)$ and $g(r)$ (Eqs~\ref{eq:fdef}--\ref{eq:gdef}), multiplying by $r$, and collecting terms of
odd power in $r$ (which become polynomial after division by $r$):
\begin{equation}
  \boxed{%
    \phi_n = \frac{(i\omega/\beta)^{2n+1}}{(2n+1)!\;4\pi\rho\beta^2}
  }
\end{equation}
\begin{equation}
  \boxed{%
    \psi_n = \frac{1}{4\pi\rho}
    \left[
      \frac{(i\omega/\alpha)^{2n+3}}{\alpha^2}
      - \frac{(i\omega/\beta)^{2n+3}}{\beta^2}
    \right]
    \frac{1}{(2n+3)!}
  }
\end{equation}
In practice, $N_{\text{Taylor}} = 8$ terms suffice for convergence at the
frequencies of interest ($\omega a/\beta < 1$).

\subsection{Smooth Second Derivatives}

The second spatial derivatives of $G_{ij}^{\mathrm{s}}$ are computed via the
chain rule, using $\partial r^2/\partial x_k = 2x_k$.  The three independent
components needed for the cube's $A$, $B$, $C$ decomposition (Eq~\ref{eq:ABCdef}) are:
\begin{align}
  G_{11,11}^{\mathrm{s}}
    &= 4x_1^2\,\Phi'' + 2\Phi' + 2\Psi + 10x_1^2\,\Psi'
       + 4x_1^4\,\Psi'',
  \\
  G_{11,22}^{\mathrm{s}}
    &= 4x_2^2\,\Phi'' + 2\Phi' + 2x_1^2\,\Psi'
       + 4x_1^2 x_2^2\,\Psi'',
  \\
  G_{12,12}^{\mathrm{s}}
    &= \Psi + 2(x_1^2 + x_2^2)\,\Psi'
       + 4x_1^2 x_2^2\,\Psi'',
\end{align}
where primes denote $d/du$ with $u = r^2$.  These are polynomial in
$(x_1, x_2, x_3)$---completely free of singularities at the
origin---and are integrated over the cube by GL quadrature.

The integrals give $A^c$, $B^c$, $C^c$ via (Eq~\ref{eq:ABCdef}):
\begin{equation}
  A^c = \int G_{11,22}^{\mathrm{s}}\,d^3x, \quad
  B^c = \int G_{12,12}^{\mathrm{s}}\,d^3x, \quad
  C^c = \int G_{11,11}^{\mathrm{s}}\,d^3x - A^c - 2B^c.
\end{equation}


%=====================================================================
\section{Displacement-Traction Basis}
%=====================================================================

For integration into the Kennett reflectivity propagator framework
(Section~20.2), the T-matrix must operate on the displacement-traction
state vector rather than the strain tensor.  This section derives the
$6\times 6$ T-matrix in the
$(u_z, u_x, u_y, t_{zz}, t_{xz}, t_{yz})$ basis.

\subsection{Coordinate System}

We adopt a right-handed coordinate system with $z$~(down), $x$~(right),
$y$~(out of page).  The mapping from the generic $(1,2,3)$ Cartesian indices
used in Parts~I--IV to $(z,x,y)$ is: $1\to z$, $2\to x$, $3\to y$.  The
Voigt ordering is:
\begin{equation}
  \boldsymbol{\varepsilon}_V
  = [\varepsilon_{zz},\;\varepsilon_{xx},\;\varepsilon_{yy},\;
     2\varepsilon_{xy},\;2\varepsilon_{zy},\;2\varepsilon_{zx}]^T.
\end{equation}

\subsection{Strain Extraction from Displacement-Traction}

For a plane wave with horizontal wavenumbers $(k_x,k_y)$, the six
independent strains can be expressed in terms of the six
displacement-traction components \emph{without requiring~$k_z$}.  The
traction components absorb the vertical-wavenumber dependence through the
stress-strain relation.  The key relations are:
\begin{align}
  \varepsilon_{zz}
    &= \frac{t_{zz} - \lambda\,ik_x\,u_x - \lambda\,ik_y\,u_y}
            {\lambda + 2\mu},
  \\
  \varepsilon_{xx} &= ik_x\,u_x, \qquad
  \varepsilon_{yy} = ik_y\,u_y,
  \\
  2\varepsilon_{zx} &= \frac{t_{xz}}{\mu}, \qquad
  2\varepsilon_{zy} = \frac{t_{yz}}{\mu},
  \\
  2\varepsilon_{xy} &= ik_y\,u_x + ik_x\,u_y.
\end{align}
These define a $6\times 6$ strain extraction matrix $\bS$ such that
$\boldsymbol{\varepsilon}_V = \bS(k_x,k_y)\;\bv^0$
where $\bv = (u_z, u_x, u_y, t_{zz}, t_{xz}, t_{yz})^T$:
\begin{equation}
\bS
= \begin{pmatrix}
  0 & -\dfrac{\lambda\,ik_x}{M} & -\dfrac{\lambda\,ik_y}{M}
    & \dfrac{1}{M} & 0 & 0 \\[8pt]
  0 & ik_x & 0 & 0 & 0 & 0 \\
  0 & 0 & ik_y & 0 & 0 & 0 \\
  0 & ik_y & ik_x & 0 & 0 & 0 \\
  0 & 0 & 0 & 0 & 0 & \dfrac{1}{\mu} \\[8pt]
  0 & 0 & 0 & 0 & \dfrac{1}{\mu} & 0
\end{pmatrix},
\end{equation}
where $M = \lambda + 2\mu$.  Note that $u_z$ does not appear explicitly:
the vertical strain $\varepsilon_{zz}$ is extracted entirely from the
traction $t_{zz}$ and the horizontal displacement components.

\subsection{Voigt T-Matrix}

The $6\times 6$ Voigt T-matrix (Eqs~\ref{eq:TVblock}--\ref{eq:Sblock}) in the $(z,x,y)$ coordinate
system is:
\begin{equation}
  \bT_V
  = \begin{pmatrix} \bD & \mathbf{0} \\
                     \mathbf{0} & \bS_{\mathrm{shear}} \end{pmatrix},
  \qquad
  \bD = T_1\,\bJ + (2T_2 + T_3)\,\bI_3,
  \qquad
  \bS_{\mathrm{shear}} = 2T_2\,\bI_3.
\end{equation}
This maps Voigt strain $\to$ Voigt effective stress perturbation:
$\Delta\boldsymbol{\sigma}^*_V
= \bT_V\,\boldsymbol{\varepsilon}^0_V$.

\subsection{Effective Stiffness Contrast Matrix}

The effective stiffness contrast tensor $\Delta c^*_{mnlp}$ in Voigt form
has cubic symmetry:
\begin{equation}
  \Delta\mathbf{c}^*_V
  = \begin{pmatrix} \bD^* & \mathbf{0} \\
                     \mathbf{0} & \bS^* \end{pmatrix},
\end{equation}
where
\begin{equation}
  D^*_{ii} = \Delta\lambda^* + 2\Delta\mu^{*,\mathrm{diag}}, \quad
  D^*_{ij} = \Delta\lambda^* \; (i \neq j), \quad
  \bS^* = 2\Delta\mu^{*,\mathrm{off}}\,\bI_3,
\end{equation}
with the effective contrasts $\Delta\lambda^*$,
$\Delta\mu^{*,\mathrm{diag}}$, $\Delta\mu^{*,\mathrm{off}}$ given by
Eqs~\eqref{eq:Dmuoff}--\eqref{eq:Dlamstar_c}.

\subsection{Traction Projection}

The traction on a $z = \text{const}$ surface is extracted from Voigt stress
by a $3\times 6$ projection matrix~$\bP$:
\begin{equation}
  \begin{pmatrix} t_{zz} \\ t_{xz} \\ t_{yz} \end{pmatrix}
  = \bP\,\boldsymbol{\sigma}_V, \qquad
  \bP
  = \begin{pmatrix}
    1 & 0 & 0 & 0 & 0 & 0 \\
    0 & 0 & 0 & 0 & 0 & 1 \\
    0 & 0 & 0 & 0 & 1 & 0
  \end{pmatrix},
\end{equation}
where $\sigma_{zx}$ is Voigt index~6 and $\sigma_{zy}$ is Voigt index~5.

\subsection{The $6\times 6$ T-Matrix in the $(\bu,\bt)$ Basis}

The full T-matrix in the displacement-traction basis maps
$(\bu,\bt)^0 \to$ scattered $(\bu,\bt)$.  Its
structure is:
\begin{equation}
  \boxed{%
    \bT_{6\times 6}
    = V\,\bigl[\text{density block} + \text{stiffness block}\bigr]
  }
\end{equation}
The \textbf{density scattering} contribution to the scattered displacement
is:
\begin{equation}
  u_i^{\mathrm{scat}}
  = V\,\omega^2\,\Delta\rho^*\,\Gamma_0\;u_i^0,
\end{equation}
giving a diagonal $3\times 3$ block in the upper-left of
$\bT_{6\times 6}$.

The \textbf{stiffness scattering} contribution to the scattered traction
is:
\begin{equation}
  \bt^{\mathrm{scat}}
  = V\;\bP\;\Delta\mathbf{c}^*_V\;\bS(k_x,k_y)\;
    \bv^0.
\end{equation}
This fills the lower $3\times 6$ block.  The complete matrix is:
\begin{equation}
  \bT_{6\times 6}
  = V
  \begin{pmatrix}
    \omega^2\,\Delta\rho^*\,\Gamma_0\,\bI_3 & \mathbf{0}_3 \\
    \multicolumn{2}{c}{%
      \bP\;\Delta\mathbf{c}^*_V\;\bS(k_x,k_y)
    }
  \end{pmatrix}.
\end{equation}
The displacement scattering from stiffness contrasts (upper-right blocks)
requires the strain-Green's tensor $H_{ijk}$ evaluated at the observation
point, which is handled by the propagation framework (FFTProp).  At the
local site, the self-consistent amplification factors (Eqs~\ref{eq:Aeoff}--\ref{eq:Aucube}) already
account for this coupling.


%=====================================================================
\section{Numerical Example (Cube --- NumPy)}
%=====================================================================

Using the same parameters as Section~10:
$V_P = 5000$~m/s, $V_S = 3000$~m/s, $\rho = 2500$~kg/m$^3$.
Contrasts: $\Delta\lambda = 2$~GPa, $\Delta\mu = 1$~GPa,
$\Delta\rho = 100$~kg/m$^3$.
Frequency $f = 10$~Hz, cube half-width $a = 10$~m.

\medskip
\begin{center}
\begin{tabular}{ll}
\toprule
\textbf{Quantity} & \textbf{Value (NumPy hybrid method)} \\
\midrule
$k_P a$ & $0.1257$ \\
$k_S a$ & $0.2094$ \\
\midrule
$\Gamma_0$ & $2.604567\mathrm{e}{-09} + 4.347360\mathrm{e}{-10}\,i$ \\
$A^c$      & $0.000000\mathrm{e}{+00} - 8.607635\mathrm{e}{-14}\,i$ \\
$B^c$      & $0.000000\mathrm{e}{+00} + 3.973958\mathrm{e}{-14}\,i$ \\
$C^c$      & $0.000000\mathrm{e}{+00} - 1.044478\mathrm{e}{-19}\,i$ \\
\midrule
$T_1^c$    & $0.000000\mathrm{e}{+00} + 2.252429\mathrm{e}{-04}\,i$ \\
$T_2^c$    & $0.000000\mathrm{e}{+00} - 4.633677\mathrm{e}{-05}\,i$ \\
$T_3^c$    & $0.000000\mathrm{e}{+00} - 2.088956\mathrm{e}{-10}\,i$ \\
\midrule
$A_u$                    & $1.001029\mathrm{e}{+00} + 1.719804\mathrm{e}{-04}\,i$ \\
$A_\theta$               & $9.999997\mathrm{e}{-01} + 5.830547\mathrm{e}{-04}\,i$ \\
$A_e^{\mathrm{off}}$     & $1.000000\mathrm{e}{+00} - 9.267354\mathrm{e}{-05}\,i$ \\
$A_e^{\mathrm{diag}}$    & $1.000000\mathrm{e}{+00} - 9.267375\mathrm{e}{-05}\,i$ \\
\midrule
$\Delta\rho^{*,c}$           & $1.001029\mathrm{e}{+02} + 1.719804\mathrm{e}{-02}\,i$ \\
$\Delta\lambda^{*,c}$        & $1.999999\mathrm{e}{+09} + 1.616595\mathrm{e}{+06}\,i$ \\
$\Delta\mu^{*,\mathrm{diag}}$ & $1.000000\mathrm{e}{+09} - 9.267375\mathrm{e}{+04}\,i$ \\
$\Delta\mu^{*,\mathrm{off}}$  & $1.000000\mathrm{e}{+09} - 9.267354\mathrm{e}{+04}\,i$ \\
\midrule
$\Delta\mu^{*,\mathrm{diag}} - \Delta\mu^{*,\mathrm{off}}$
  & $-3.886223\mathrm{e}{-05} - 2.088956\mathrm{e}{-01}\,i$ \\
\bottomrule
\end{tabular}
\end{center}

\medskip
The values of $T_1^c$ and $T_2^c$ agree with the SymPy symbolic computation
to approximately $10^{-8}$ relative error.  $T_3^c$ and $C^c$ differ by
approximately $0.3\%$, because the NumPy implementation retains higher-order
Taylor terms (8~terms \textit{vs.}\ the $O((\omega a)^7)$ truncation in the
symbolic code).

\bigskip
\noindent
\textit{Generated from \texttt{CubicTMatrix/effective\_contrasts.py} and
\texttt{CubicTMatrix/voigt\_tmatrix.py} using NumPy hybrid computation.
$\Gamma_0$ evaluated by 3D Gauss--Legendre quadrature ($32^3$~points);
$A^c$, $B^c$, $C^c$ evaluated by polynomial Taylor expansion (8~terms)
with GL quadrature.}


% --- Bibliography ------------------------------------------------------------
\bibliographystyle{gji}
\bibliography{references}

\end{document}
